\section{Scope, Definitions, and Boundaries}
\label{sec:scope-definitions}

The central object of this survey is a \emph{course}: an ordered spatial task
specification that determines where a robot is to travel and how traversal is
recognized as progress or completion. ``Ordered'' is substantive. A course is
not merely a collection of geometric primitives, obstacles, or visual assets; it
encodes an intended traversal relation among task-defining elements. A ground
course can therefore be a closed or open corridor with an associated route order,
an aerial course a sequence of gate poses, and a maritime course a sequence of
waypoints, markers, or buoy-defined gates. These distinctions preserve the
semantics that make geometry consequential to the task: a gate passed from the
wrong side, or a buoy encountered out of order, need not constitute completion.
Ground racing environments, drone-racing tracks with gates, and maritime
competition tasks illustrate these different realizations of the same abstraction
\cite{KlimovNodateCarRacing,Kaufmann2018DeepDrone,RoboNation2026RoboBoat2026}.

The definition is intentionally independent of a particular file format, simulator,
or vehicle. It includes the spatial constraints and traversal semantics needed to
state the task, but does not imply a unique dynamics model, reward, sensor suite,
or controller. A course may be closed, open, branching, or include prescribed
choices, provided that its completion semantics are declared. It may also carry
task-relevant spatial attributes such as corridor width, gate orientation, clearance,
or permitted passage direction. The normative interface by which those attributes
are serialized and tested is specified separately in \cref{sec:benchmark-protocol};
here the purpose is to establish distinctions that prevent incompatible objects from
being treated as interchangeable evidence.

\subsection{Objects and transformations}

We distinguish the following objects and operations.

\paragraph{Course object.}
The \emph{course object} is the abstract task-defining spatial entity: a traversable
corridor and route relation, an ordered gate chain, a waypoint sequence, or a
buoy course. It specifies what is to be traversed, not how it is stored or rendered.
For a road network, a graph alone is insufficient unless the relevant route or
completion relation is also stated; for an aerial gate course, positions alone are
insufficient when gate orientation or crossing direction matters.

\paragraph{Representation.}
A \emph{representation} is an encoding of a course object. Examples include a
centerline with width profile, paired boundaries, a segment grammar, a waypoint
graph, an ordered list of gate poses, or a simulator-native world description.
Representations may be lossless or deliberately task-specific projections. They
are interfaces between generation, validation, simulation, and evaluation, rather
than generation methods in themselves. This distinction is material because
benchmark and standard artifacts can prescribe composable scenarios, route
semantics, or exchange formats without constructing new courses
\cite{Althoff2017CommonRoadComposable,ASAM2024ASAMOpenSCENARIO}.

\paragraph{Concrete realization or artifact.}
A \emph{realization} is a specific instantiated artifact that implements a
representation: for example, a map, world file, mesh, occupancy image, track asset,
or scenario record. The realization may add collision geometry, visual markers,
textures, spawn locations, and simulator metadata. Such additions are relevant to
interoperability, but they do not retroactively change the course object unless they
alter the task-defining traversal constraints. A fixed map collection or competition
layout is consequently evidence about representation, interface, or constraint
requirements, not evidence of a generation method.

\paragraph{Course generation and artifact management.}
\emph{Course-generation operations} create or change task-defining spatial
structure: synthesis creates a new course, mutation changes a candidate, and repair
projects a candidate toward declared constraints. \emph{Artifact-management
operations} instead select courses from an extant set or replay previously
constructed artifacts, thereby managing course populations without generating new
course structure. Serialization changes a course's encoding or concrete realization.
It is covered because portable realizations matter to interoperability, but
serialization alone is not course generation. These operations are not mutually
exclusive method families: one pipeline can synthesize candidates, reject or repair
them, select a subset for training, and serialize retained artifacts. Procedural road
construction, replay-guided level selection, and diversity-guided search illustrate
distinct operations on course structure and course populations
\cite{Li2020ImprovingGeneralization,Jiang2021ReplayGuided,Blattner2024DiversityGuided}.

\paragraph{Validity predicate.}
A \emph{validity predicate} is a declared test on a course, its representation, or
its realization. It may concern geometry (for example, clearance or curvature),
topology (closure, connectivity, or non-self-intersection), ordered-task semantics,
dynamics-conditioned traversability, or successful import into a simulator. Validity
is orthogonal to generation: rejection, construction by constraint, repair, and
post-hoc simulation checks are alternative ways to relate a generator to one or
more predicates. It is therefore incorrect to classify a validity test as a
generation method, or to infer that a syntactically valid artifact is dynamically
traversable. Search-based test generation makes this separation explicit by
reporting invalid candidates alongside the exploration procedure
\cite{Blattner2024DiversityGuided}.

\paragraph{Training distribution and evaluation suite.}
A \emph{training distribution} is a population of course artifacts together with
the sampling rule used while fitting, tuning, or adapting a controller or planner.
It can be fixed, procedurally sampled, curated by selection, or changed over time.
An \emph{evaluation suite} is instead a fixed or reproducibly sampled, versioned
set of course instances and conditions used to measure a system. The same generator
can define a training distribution without defining an evaluation suite, and a fixed
benchmark can define an evaluation suite without being a generator. This separation
is essential when procedural environments are used to study generalization
\cite{Cobbe2020LeveragingProcedural}.

\paragraph{Controller and planner.}
A \emph{controller} or \emph{planner} maps observations, state estimates, and goals
to actions or trajectories on a given course. Racing-line optimization, trajectory
tracking, and progress rewards operate \emph{within} a course instance. They may
expose constraints that a generator should satisfy, but they are not course
generation. A control or planning study contributes to this survey's scope only when
it additionally creates or changes course structure, or selects or replays course
artifacts as an explicit artifact-management contribution.

\begin{table}[t]
  \centering
  \caption{Boundary of surveyed contributions and evidence roles. Classification
  follows the role established for course structure or artifacts, rather than the
  application domain in which it appears.}
  \label{tab:scope-boundaries}
  \small
  \begin{tabularx}{\linewidth}{@{}>{\raggedright\arraybackslash}p{0.24\linewidth}Y Y@{}}
    \toprule
    Area & In-scope contribution or evidence role & Out-of-scope use or inference \\
    \midrule
    Course generation & Synthesis, mutation, or repair of a course-defining route,
    corridor, gate, waypoint, buoy, or world geometry and its traversal semantics &
    Racing actions, a trajectory on a fixed course, or a vehicle model \\
    Artifact management and interoperability & Selection or replay of course
    populations; serialization of a representation or realization & An inference
    that selection, replay, or serialization alone generates course structure \\
    Validity and repair & A declared validity predicate as evidence; repair as a
    transformation of a candidate course artifact & A claim that all generated
    artifacts are valid without a stated predicate and population \\
    Benchmark, rule, or standard & Representation, interface, fixed suite, or
    admissible task constraint & A method merely because it supplies a fixed layout
    or competition rule \\
    PCG and randomization & Spatial task structure that defines the course & Texture,
    lighting, weather, sensor noise, traffic, or dynamics on an unchanged course \\
    Control and planning & An explicit additional contribution that synthesizes,
    mutates, or repairs course structure, or selects or replays course artifacts & A
    racing line, policy, or optimizer evaluated only on pre-existing artifacts \\
    \bottomrule
  \end{tabularx}
\end{table}

\subsection{Boundary cases}

\Cref{tab:scope-boundaries} establishes the survey boundary by contribution and
evidence role. Generic procedural content generation (PCG) is in scope only to the
extent that it generates course structure or manages the spatial artifacts that
define the robot's traversal task. Environment
randomization is not course generation when it varies appearance, sensing, dynamics,
or traffic while leaving the course object unchanged. Conversely, a method that
changes a road graph, corridor, gate sequence, or buoy layout can be relevant even if
its original application is testing, simulation, or a non-racing navigation domain.
The criterion is transferability of the represented task structure, not a shared
application label.

Fixed benchmarks, competition rules, and standards occupy a distinct supporting
role. They may establish coordinate conventions, ordering requirements, route or
gate semantics, permissible geometry, serialization interfaces, and evaluation
conditions. F1TENTH and CommonRoad, for example, provide evaluation artifacts or
scenario interfaces; they should not be counted as generators simply because a
controller is evaluated on their courses \cite{OKelly2020F1TENTHOpen,Althoff2017CommonRoadComposable}.
Likewise, official maritime rules can define buoy-course completion and scoring
without supplying an artifact-synthesis procedure
\cite{RoboNation20222022Maritime,RoboNation2026RoboBoat2026}.

These distinctions also prevent an invalid aggregation. A retained source can
contribute a representation, a fixed-suite interface, a validity predicate, a
course-generation operation, an artifact-management operation, an interoperability
role, or a reporting requirement; it need not contribute a distinct generation
method. The survey therefore treats source roles as multi-label evidence rather than
treating source count as method count. Quantitative conclusions require completed,
reliable coding and an explicit unit of analysis.
